\newcommand{\tipoRequerimento}{Contestação de Indeferimento por Ausência de Prova Técnica da Carga}

% Contestação para o caso de a concessionária não apresentar provas que fundamentem o indeferimento da reclamação, DEVE SER ANÁLISADO DE QUALQUER MANEIRA"

\section{DOS FATOS E DOS DIREITOS}

\begin{enumerate}[resume]
    \item Em resposta ao requerimento de reclassificação tarifária da Unidade Consumidora \unidadeConsumidora{}, a concessionária \nomeConcessionaria{} indeferiu o pleito alegando, de forma genérica e superficial, que a referida unidade não possui circuito exclusivo dedicado à Iluminação Pública.

    \item Para fundamentar a sua recusa, a distribuidora limitou-se a anexar capturas de tela (\textit{prints}) extraídas de redes sociais relativas a eventos comunitários e escolares ocorridos no local, arguindo presuntivamente que a presença de praça, quadra ou espaço de eventos descaracterizaria a exclusividade de carga do sistema de iluminação pública.

    \item Ocorre que registros fotográficos de eventos sociais ou comunitários não constituem prova técnica hábil nem instrumento regulatório idôneo para atestar a existência ou a operação permanente de cargas elétricas alheias ao sistema de iluminação pública no ponto de medição.

    \item Nos termos do art. 189, inciso II, da Resolução Normativa ANEEL nº 1.000/2021, enquadram-se na classe Iluminação Pública os bens públicos destinados ao uso comum do povo, tais como praças, parques, jardins e quadras esportivas, ainda que sujeitos a restrições de horário, cercamento ou eventos promovidos pela administração municipal.

    \item A desqualificação do enquadramento na classe de Iluminação Pública não pode ser calcada em meras presunções ou conjecturas visuais da concessionária. Para infirmar a exclusividade do circuito de IP, incumbe à distribuidora o ônus regulatório e probatório de instruir a sua decisão com elementos técnicos objetivos e conclusivos.

    \item A idoneidade do indeferimento exige a apresentação de um Termo/Relatório de Vistoria Técnica de Campo instruído com:
    \begin{enumerate}
        \item Inventário discriminado de toda a Carga Instalada vinculada à unidade consumidora;
        \item Registro fotográfico do medidor, do quadro de distribuição e dos circuitos derivados;
        \item Identificação clara e mensurável de eventuais tomadas ou equipamentos de consumo permanente que descaracterizem a exclusividade do circuito de iluminação;
        \item Identificação do responsável técnico pela inspeção.
    \end{enumerate}

    \item A omissão desses elementos demonstra a patente fragilidade e a superficialidade do indeferimento administrativo exarado pela concessionária, caracterizando manifesta violação ao dever de motivação, à transparência e aos princípios regulatórios do setor elétrico.

    \item Diante da ausência de laudo de carga ou de qualquer comprovação técnica de inexistência de exclusividade no circuito, a recusa da \nomeConcessionaria{} em reclassificar a UC \unidadeConsumidora{} gera faturamento indevido continuado, ensejando a devolução em dobro dos valores pagos a maior nos últimos 10 (dez) anos, conforme o art. 323 da REN nº 1.000/2021, o Despacho ANEEL nº 2.006/2024 e o art. 205 do Código Civil.
\end{enumerate}

\section{DOS PEDIDOS}

\begin{enumerate}[resume]
\item Diante do exposto, requer que a Concessionária \nomeConcessionaria{}:

\begin{enumerate}
    \item Proceda com o reexame do indeferimento e efetive a reclassificação tarifária da UC \unidadeConsumidora{} para a classe Iluminação Pública (subgrupo B4a), ante a absoluta ausência de laudo de vistoria técnica que demonstre a não exclusividade do circuito de IP;

    \item Subsidiariamente, caso insista na tese de não exclusividade da carga, que apresente o Relatório/Termo de Vistoria de Campo completo contendo o inventário discriminado e fundamentado de toda a carga instalada e dos circuitos elétricos da referida UC;

    \item Caso a concessionária não possua o inventário discriminado de carga, que agende nova vistoria técnica conjunta no local, a fim de realizar a conferência pormenorizada dos circuitos \textit{in loco};

    \item Efetue a devolução em dobro de todos os valores cobrados a maior em razão da tarifa inadequada aplicada à unidade consumidora, abrangendo o período prescricional de 10 (dez) anos, devidamente atualizados pelo IPCA e acrescidos de juros de mora de 1\% ao mês calculados \textit{pro rata die}, nos termos da REN ANEEL nº 1.000/2021 e do Despacho ANEEL nº 2.006/2024;

    \item Realize o ressarcimento dos valores apurados por meio de depósito bancário em conta corrente de titularidade do Município, conforme estabelecido no § 7º do art. 323 da REN ANEEL nº 1.000/2021.
\end{enumerate}
\end{enumerate}